\begin{figure}[t]
    \centering
    \includegraphics[width=0.9\textwidth]{sampling.pdf}
    \caption{Illustration of different configuration sampling methods for exploring the potential energy landscapes in atomistic modelling.}
    \label{fig:sampling_methods}
\end{figure}
In \cref{eq:general_sampling_1} and \cref{eq:general_sampling_2},  we have established the statistical viewpoint on the physical/chemical properties for a given material system as the ensemble average of the corresponding property for individual configuration on the PES, which is drawn from seven different scenarios highlighted in \cref{sect:sampling_examples}  where sampling the PES are relevant in the fundamental materials science. In order to perform the computation as in \cref{eq:general_sampling_1} or \cref{eq:general_sampling_2}, the  first piece of recipe is to have a suitable energy model with which we can compute the corresponding energy of a given material with known atomistic structure $\bm{R}$, that we have  reviewed comprehensively in \cref{chapter:energy_models}. However, very often, we may not have a full dataset of possible structures for a given material readily available, in which case, we need a method to pragmatically sample the PES to generate a distribution of stable structures $\mathcal{P}(E(\bm{R}))$ at a given thermodynamic ensemble $(P,V,T,\mu)$ specified by the thermodynamic variables of pressure $T$, volume $V$, temperature $T$ and chemical potential $\mu$. Having access to these structures is the pre-requisite before subsequent property calculations can be conducted, thus PES sampling constitutes a critical component in modern computational materials science, which will be the focus for this section. More specifically, we will be focusing on the \textit{traditional} PES sampling methods as illustrated in \cref{fig:sampling_methods}, which exclude those driven by the recently popular generative AI models, that will be reserved for the later sections.

However, we do like to highlight here, that both the PES sampling and generative AI methods are targeted towards conquering the probability distribution $P(E(\bm{R}))$, which is  the Boltzmann distribution 
\begin{equation}
    \label{eq:Boltzmann_distr}
    P(E)=e^{-\beta E}/\mathcal{Z},
\end{equation}
in which $\beta=1/k_B T$, and $\mathcal{Z}$ is the normalisation constant as given in \cref{eq:configuration_partiton_function}. This seemingly simple distribution is generally unknown for a given material system because the normalisation constant $\mathcal{Z}$ is determined based on the energy of all configurations on the PES. As the potential energy function $E(\bm{R})$ lies in a high-dimensional space with complexity that grows with the dimension of $\bm{R}$, the number of potential local minima becomes  analytically intractable, leading to the difficulty of obtaining $\mathcal{Z}$ deterministically. In this sense, methods for sampling the PES are developed, primarily, in the hope to find:
\begin{enumerate}
    \item the global minimum, which corresponds to the most stable configuration on the PES, and
    \item  a set of low-energy metastable configurations that contribute significantly towards the Boltzmann weight $e^{-\beta E}$.
\end{enumerate}

Broadly speaking, methods for PES sampling may be classified into two classes of approaches, namely, the ones with energy guidance and the ones without. The well-known Monte Carlo [MC, \cref{fig:sampling_methods}(a)] and molecular dynamics [MD, \cref{fig:sampling_methods}(b)] are typical examples of energy-guided sampling approaches, whereby the evolution of the configurations and/or the acceptance of new configurations are tightly coupled with the underlying energy model employed. On the one hand, this allows us to directly generate the corresponding Boltzmann distribution at the sampling temperature [\cref{eq:Boltzmann_distr}]. On the other hand, it can encounter difficulties or even impossibility in reaching deep minima separated by high energy barriers, such that ergodicity can never be reached.\cite{wales1997global,doye1999evolution} The energy-guided samplings will be the major focus of this section. 

In contrary, random search [\cref{fig:sampling_methods}(c)] first generates a set of candidate structures according to a specific set of system-dependent geometric rules, such as space group symmetries and molecular orientations in molecular crystals,\cite{case2016convergence}, after which energy minimisations are performed on each candidate structure to their respective local minimum. It provides us with the possibility to sample the PES globally, and the possibility to benchmark the performances of different energy models on an identical set of structures.\cite{Hunnisett:aw5094} However, one caveat behind the suitability of using random search for PES sampling is that the degrees of freedom being sampled should be as independent from each other as possible.\cite{echenique2006definition} In the case when multiple different degrees of freedom are coupled together, such as the torsions and bond angles in a molecule, they must be changed cooperatively in order to make a stable molecular conformation by avoiding steric clashes, the chances of locating low-energy conformers are significantly lowered when these degrees of freedom are searched independently, rending the necessity to guide the search with the topologies of the underlying PES, such as the low-mode search.\cite{kolossvary1996low}